\textbf{Mezivrstvová polarizace} vzniká v nehomogenním dielektriku hromaděním pohyblivých nábojů na rozhraních oblastí s odlišnými elektrickými vlastnostmi. Přiložené pole oddělí náboje, které se na nehomogenitách zachytí a vytvářejí prostorové dipóly. Jde o pomalý polarizační mechanismus; oddělení nábojů znázorňuje schéma.
